\subsection*{IV 20.3: Section hyperplane générique---irréductibilité et connexité géomét. (cote 25, sites 15--18)}
The Bertini--Zariski discussion turns from the local properties of no.\ 2 to geometric irreducibility and connectedness of the generic section. It writes the function field of the incidence section in affine coordinates, reducing the argument to the lemma attributed to Zariski on site 17; the explicit birational description is itself proposed as a corollary to 2.1. The handwritten revision of that lemma replaces the original algebraicity qualification by a transcendence-degree qualification and deletes the claim equating radicial multiplicities, while retaining the assertion about primary extensions. The surrounding argument uses the generic incidence construction of no.\ 1 and the separability discussion of no.\ 2. The connectedness corollary invokes paragraph 6 and the theorem named ``Hartshorme'' in the typescript, joining irreducible components through their intersections. Its script family of closed subsets remains typographically doubtful in the checked text. Several changes matter to the reading: the theorem on site 15 loses its geometrically integral alternative; a long canceled proof continues through the first part of site 16; the new proof begins at ``Soit K/k'' there; site 17 adds the attribution ``(Zariski)'', modifies the immersion hypothesis of corollary 3.2, and inserts algebraic closedness in corollary 3.3. The author leaves Dieudonné the choice of a proof to copy from the literature. The qualification above the canceled immersion wording still has an uncertain initial group. The guide preserves the written argument and its omissions: the canceled formula on site 16 has a blank in $k(\ )$, and site 18 prints $Y_\eta=X_\eta$ while completing the connectedness proof. Sources: \href{https://grothendieck.umontpellier.fr/web/viewer.html?file=../25.pdf\#page=16}{site 15}, \href{https://grothendieck.umontpellier.fr/web/viewer.html?file=../25.pdf\#page=17}{site 16}, \href{https://grothendieck.umontpellier.fr/web/viewer.html?file=../25.pdf\#page=18}{site 17}, and \href{https://grothendieck.umontpellier.fr/web/viewer.html?file=../25.pdf\#page=19}{site 18}; the paragraph-6 reference is the one given in the prenotes.
